%% kapa: exam=1 round=27 year=2016 period=1 form=A subject=civil_law subject_ko=민법 seq=38 number=38
%% source: Q-Net 감정평가사 기출문제 (cst00309) artlSeq=5169402 file=extracted/1008/5169402/제27회 감정평가사 1차 문제/제27회 감정평가사 1교시문제(A형).pdf
%% conversion: kordoc 4.24.51 -> tools/gen1.py (auto); review=proofread-sonnet; flags=none
\begin{question}{38}
지역권에 관한 설명으로 옳지 않은 것은?
\begin{choices}
  \item 민법상 지역권의 존속기간은 최장 30년이지만 갱신할 수 있고, 이를 등기하여 제3자에 대항할 수 있다.
  \item 요역지와 승역지는 반드시 서로 인접할 필요가 없다.
  \item 공유자의 1인이 지역권을 취득하는 때에는 다른 공유자도 이를 취득한다.
  \item 지역권설정등기는 승역지의 등기부 을구에 기재된다.
  \item 지역권자는 지역권을 방해할 염려있는 행위를 하는 자에 대하여 그 예방을 청구할 수 있다.
\end{choices}
\end{question}
\answer{1}
